\chapter{Unit 1: Foundations of Learning, Creativity and Design Thinking}
\label{chap:unit1_qb}

\section*{Practice Question Set (30 MCQs)}
\addcontentsline{toc}{section}{Practice Question Set (30 MCQs)}

\mcqitem{1}{What is the fundamental starting point of Design Thinking?}{Cutting-edge technology}{Market profit maximization}{Human needs and user empathy}{Government regulations}

\mcqitem{2}{Who coined the term "Wicked Problems" in 1973 to describe complex, ill-defined social challenges?}{Herbert Simon and Peter Rowe}{Horst Rittel and Melvin Webber}{David Kelley and Tim Brown}{Karl Duncker and David Kolb}

\mcqitem{3}{Which of the following is NOT one of IDEO's three overlapping lenses of innovation?}{Desirability}{Feasibility}{Viability}{Predictability}

\mcqitem{4}{What type of thinking focuses on generating a broad variety of new and novel ideas?}{Convergent thinking}{Divergent thinking}{Critical thinking}{Deductive reasoning}

\mcqitem{5}{In cognitive learning, why is active retrieval superior to passive re-reading?}{It takes less time}{It strengthens neural memory traces and recall pathways}{It avoids critical analysis}{It relies exclusively on short-term memory}

\mcqitem{6}{Which brain mode is characterized by relaxed attention and subconscious incubation of ideas?}{Focused mode}{Diffuse mode}{Analytical mode}{Reflexive mode}

\mcqitem{7}{What is "Functional Fixedness"?}{Inability to change hardware components}{A cognitive bias where an object is seen only in terms of its traditional use}{The capacity to repair broken tools}{A systematic software testing methodology}

\mcqitem{8}{Karl Duncker's famous 1945 experiment demonstrating functional fixedness is known as the:}{Box and String Test}{Candle Problem}{Nine Dots Puzzle}{Tower of Hanoi}

\mcqitem{9}{Which sequence accurately represents the four stages of Kolb's Experiential Learning Cycle?}{Remember $\rightarrow$ Understand $\rightarrow$ Apply $\rightarrow$ Analyze}{Concrete Experience $\rightarrow$ Reflective Observation $\rightarrow$ Abstract Conceptualization $\rightarrow$ Active Experimentation}{Ideate $\rightarrow$ Define $\rightarrow$ Prototype $\rightarrow$ Test}{Plan $\rightarrow$ Do $\rightarrow$ Check $\rightarrow$ Act}

\mcqitem{10}{What does the acronym VARK stand for in learning styles?}{Verbal, Analytical, Reasoning, Knowledge}{Visual, Aural, Read/Write, Kinesthetic}{Vision, Auditory, Reality, Kinetics}{Variable, Action, Reflex, Kinetic}

\mcqitem{11}{What is the highest cognitive level in Bloom's Revised Taxonomy?}{Analyze}{Evaluate}{Create}{Apply}

\mcqitem{12}{What distinguishes an "Innovation" from an "Invention"?}{Innovation does not require technology}{Innovation involves successful implementation and real-world adoption of value}{Inventions are always cheaper than innovations}{Innovation is always patented}

\mcqitem{13}{In which book did Herbert Simon state that design is "transforming existing conditions into preferred ones"?}{The Sciences of the Artificial (1969)}{Design Thinking (1987)}{The Design of Everyday Things (1988)}{Change by Design (2009)}

\mcqitem{14}{The German design school established in 1919 that united art, craft, and functional design was:}{Bauhaus}{Stanford d.school}{Ulm School}{Royal College of Art}

\mcqitem{15}{Who authored the 1987 book titled *Design Thinking*, focusing on architectural reasoning?}{Tim Brown}{Peter Rowe}{Don Norman}{Alex Osborn}

\mcqitem{16}{In which year was the Stanford d.school (Hasso Plattner Institute of Design) founded?}{1991}{1999}{2005}{2015}

\mcqitem{17}{What are the 5 iterative modes of the Stanford d.school Design Thinking process?}{Listen, Think, Build, Deploy, Maintain}{Empathize, Define, Ideate, Prototype, Test}{Discover, Design, Develop, Deliver, Debrief}{Plan, Draft, Review, Publish, Archive}

\mcqitem{18}{What are the three phases of the IDEO Human-Centred Design framework?}{Input, Process, Output}{Inspiration, Ideation, Implementation}{Investigation, Iteration, Integration}{Initiation, Incubation, Impact}

\mcqitem{19}{Which product design success featured thick, comfortable handles originally developed for arthritis sufferers?}{Juicero}{OXO Good Grips}{Apple iPhone}{Ford Edsel}

\mcqitem{20}{What made Kenya's M-Pesa a global design thinking success in 2007?}{High-end 4G smartphone apps}{Simple SMS-based money transfers supported by local agent networks}{Zero-fee luxury banking lounges}{Crypto-currency smart contracts}

\mcqitem{21}{Why is the Jaipur Foot considered a human-centred design triumph in India?}{It uses carbon fiber imported from Germany}{It is customized for local customs like barefoot walking, farming, and squatting at low cost}{It includes built-in electronic tracking}{It was designed without consulting disabled users}

\mcqitem{22}{Why did the Juicero press (2017) fail commercially?}{It was prone to electrical fires}{Users discovered the expensive proprietary juice bags could easily be squeezed by hand}{It lacked Wi-Fi connectivity}{It was banned by FDA regulations}

\mcqitem{23}{Why did the consumer edition of Google Glass face severe market backlash?}{Poor battery life only}{Privacy concerns regarding hidden recording, social awkwardness, and unclear everyday value}{High weight causing physical injury}{Inability to connect to the internet}

\mcqitem{24}{Deliberate practice is characterized by which of the following?}{Mindless repetition without feedback}{Focused effort on weaknesses combined with immediate corrective feedback}{Multitasking during learning}{Memorizing answers without understanding}

\mcqitem{25}{Metacognition is best defined as:}{The speed of brain processing}{Thinking about and regulating one's own learning and thinking processes}{Emotional intelligence in social groups}{Telepathic connection between team members}

\mcqitem{26}{What is the capacity of the brain to reorganize itself by forming new neural connections throughout life?}{Neurogenesis}{Neuroplasticity}{Functional fixedness}{Synaptic pruning}

\mcqitem{27}{In creative problem solving, what role does a failed prototype play?}{Total loss of project budget}{Cheap and rapid discovery of flawed assumptions}{Proof that the design team lacks talent}{Reason to permanently abandon the project}

\mcqitem{28}{Which thinking mode applies logic, filters criteria, and evaluates constraints to choose the best solution?}{Divergent thinking}{Convergent thinking}{Lateral thinking}{Associative thinking}

\mcqitem{29}{What is a "Mental Set"?}{A positive mindset before starting work}{The tendency to rely rigidly on previously successful methods even when simpler approaches exist}{A collection of brainstorming tools}{Psychological readiness to accept feedback}

\mcqitem{30}{Which of the following is a primary countermeasure to Functional Fixedness?}{Sticking to the user manual}{Reframing the problem and deconstructing objects into fundamental properties}{Increasing project deadlines}{Asking for managerial approval}


\newpage
\section*{Answer Key \& Step-by-Step Explanations}
\addcontentsline{toc}{section}{Answer Key \& Explanations}

\begin{answerkeytable}{Unit 1: Quick Answer Key}
\begin{tabular}{c c c c c c c c c c}
\toprule
\textbf{Q1} & \textbf{Q2} & \textbf{Q3} & \textbf{Q4} & \textbf{Q5} & \textbf{Q6} & \textbf{Q7} & \textbf{Q8} & \textbf{Q9} & \textbf{Q10} \\
\textbf{C} & \textbf{B} & \textbf{D} & \textbf{B} & \textbf{B} & \textbf{B} & \textbf{B} & \textbf{B} & \textbf{B} & \textbf{B} \\
\midrule
\textbf{Q11} & \textbf{Q12} & \textbf{Q13} & \textbf{Q14} & \textbf{Q15} & \textbf{Q16} & \textbf{Q17} & \textbf{Q18} & \textbf{Q19} & \textbf{Q20} \\
\textbf{C} & \textbf{B} & \textbf{A} & \textbf{A} & \textbf{B} & \textbf{C} & \textbf{B} & \textbf{B} & \textbf{B} & \textbf{B} \\
\midrule
\textbf{Q21} & \textbf{Q22} & \textbf{Q23} & \textbf{Q24} & \textbf{Q25} & \textbf{Q26} & \textbf{Q27} & \textbf{Q28} & \textbf{Q29} & \textbf{Q30} \\
\textbf{B} & \textbf{B} & \textbf{B} & \textbf{B} & \textbf{B} & \textbf{B} & \textbf{B} & \textbf{B} & \textbf{B} & \textbf{B} \\
\bottomrule
\end{tabular}
\end{answerkeytable}

\begin{expbox}[Detailed Pedagogical Explanations]
\begin{enumerate}[leftmargin=6mm, itemsep=2.5mm]
  \item \textbf{[Q1: Option C]} Design Thinking always begins with user empathy and understanding human needs.
  \item \textbf{[Q2: Option B]} Horst Rittel and Melvin Webber coined "Wicked Problems" in 1973.
  \item \textbf{[Q3: Option D]} The three lenses are Desirability (people), Feasibility (technology), and Viability (business).
  \item \textbf{[Q4: Option B]} Divergent thinking expands the solution space by generating numerous diverse ideas.
  \item \textbf{[Q5: Option B]} Active recall actively engages neural memory networks, making retrieval durable.
  \item \textbf{[Q6: Option B]} Diffuse mode allows unconscious incubation and wide-ranging connections across the brain.
  \item \textbf{[Q7: Option B]} Functional fixedness limits a person to seeing an object only in its traditional role.
  \item \textbf{[Q8: Option B]} Duncker's candle problem tested using a matchbox as a shelf rather than just a container.
  \item \textbf{[Q9: Option B]} Kolb's cycle moves: Concrete Experience $\rightarrow$ Reflective Observation $\rightarrow$ Abstract Conceptualization $\rightarrow$ Active Experimentation.
  \item \textbf{[Q10: Option B]} VARK represents Visual, Aural, Read/Write, and Kinesthetic learning modalities.
  \item \textbf{[Q11: Option C]} In Bloom's revised taxonomy, 'Create' is at the very apex of cognitive skills.
  \item \textbf{[Q12: Option B]} An invention is a novel device/process; an innovation is successfully implemented to deliver value.
  \item \textbf{[Q13: Option A]} Herbert Simon wrote *The Sciences of the Artificial* in 1969.
  \item \textbf{[Q14: Option A]} Bauhaus was founded in 1919 in Germany, integrating craft, art, and industrial design.
  \item \textbf{[Q15: Option B]} Peter Rowe published *Design Thinking* in 1987.
  \item \textbf{[Q16: Option C]} The Stanford d.school was founded in 2005 by David Kelley.
  \item \textbf{[Q17: Option B]} Stanford d.school's 5 stages: Empathize, Define, Ideate, Prototype, and Test.
  \item \textbf{[Q18: Option B]} IDEO's 3I framework consists of Inspiration, Ideation, and Implementation.
  \item \textbf{[Q19: Option B]} OXO Good Grips created ergonomic kitchen tools that became a classic universal design hit.
  \item \textbf{[Q20: Option B]} M-Pesa utilized basic SMS technology and local agents to serve unbanked populations.
  \item \textbf{[Q21: Option B]} Jaipur Foot adapted prosthetics to Indian cultural practices (barefoot walking, squatting).
  \item \textbf{[Q22: Option B]} Juicero was over-engineered and solved a non-problem since juice packs were hand-squeezeable.
  \item \textbf{[Q23: Option B]} Google Glass faced severe social resistance due to perceived privacy violations and lack of clear consumer utility.
  \item \textbf{[Q24: Option B]} Deliberate practice requires focused concentration on errors with immediate feedback.
  \item \textbf{[Q25: Option B]} Metacognition is conscious awareness and control over one's own thinking.
  \item \textbf{[Q26: Option B]} Neuroplasticity is the brain's ability to rewire its neural pathways through experience and learning.
  \item \textbf{[Q27: Option B]} In design thinking, early prototype failure is an inexpensive source of validated learning.
  \item \textbf{[Q28: Option B]} Convergent thinking filters, evaluates, and selects from generated possibilities.
  \item \textbf{[Q29: Option B]} A mental set is a cognitive habit of using familiar problem-solving methods rigidly.
  \item \textbf{[Q30: Option B]} Reframing and breaking objects down into raw properties breaks functional fixedness.
\end{enumerate}
\end{expbox}
